\section{Выводы}

В этой работе был реализован словарь на основе дерева PATRICIA, поддерживающий операции вставки, удаления, поиска, сохранения и загрузки данных.

Наиболее важной частью реализации оказались операции, изменяющие внутреннюю структуру дерева: удаление вершин и восстановление дерева после чтения из файла. По сравнению с обычными линейными структурами данных, здесь особенно важно сохранять внутренние инварианты и аккуратно работать с указателями.

В результате была получена программа, удовлетворяющая требованиям задания: словарь корректно работает с регистронезависимыми строковыми ключами, поддерживает бинарное сохранение и загрузку и восстанавливает структуру дерева в идентичном виде.

Практический вывод по лабораторной работе состоит в том, что реализация нетривиальной структуры данных требует внимания не только к основным операциям поиска и вставки, но и к граничным случаям, связанным с модификацией дерева и внешним представлением данных.
\pagebreak
